\setcounter{section}{4}

\Needspace{5\baselineskip}
\hypertarget{multi-head-latent-attention}{%
\section{Multi-Head Latent Attention}\label{multi-head-latent-attention}}

\Needspace{5\baselineskip}
\hypertarget{ux4e00ux57faux672cux601dux60f3-1}{%
\subsection{一、基本思想}\label{ux4e00ux57faux672cux601dux60f3-1}}

DeepSeek提出的MLA（Multi-Head Latent Attention）旨在解决现代GPU显存带宽不足的情况下，KV Cache显存占用过大的问题。它通过低秩矩阵分解和矩阵吸收技术，在推理时极大地压缩了KV Cache的体积，又不会出现MQA、GQA的损失问题。

在线性代数中，矩阵的``秩''代表了矩阵中真正独立的、不相关的信息量。假设我们有一个 \(100\times100\) 的矩阵，如果第2行是第1行的2倍，第3行是第1行的3倍\ldots\ldots 那么虽然它100 行，但实际上只有1行是有用的信息。它的秩就是1。

在标准的 Multi-Head Attention 中，我们有 \(H\) 个头（比如32个）。每个头分开反向传播，模型会分别为每个头学习不同的投影矩阵 \(W_K\)、\(W_V\)。事实上，这32个头学到的 Key 和 Value不是完全独立的，它们之间存在极强的相关性。比如处理单词``苹果''：头1关注它的``颜色（红）''，头2关注它的``形状（圆）''，头3关注它的``类别（水果）''，这些特征虽然不同，但它们都源自同一个语义核心``苹果''。\(W_K\) 和 \(W_V\) 矩阵用于特征重组，它们事实上关注的是：头关注的特征向量可以看作原特征向量各个维度怎样的加权，既然特征之间存在强相关性，这些权重向量也就存在强相关性。因而DeepSeek MLA 的核心假设是：多头注意力（MHA）中的 Key-Value 矩阵虽然维度很高，但在数学上是``低秩''（Low-Rank）的。即：大部分参数都是冗余的。

数学上，如果一个大矩阵 \(W_{\mathrm{full}}\) 是低秩的，它就可以被无损（或极低损耗）地分解为两个低秩矩阵 \(W_{\mathrm{down}}W_{\mathrm{up}}\) 的乘积，在MLA中这两个矩阵分别被称为降维矩阵和升维矩阵。从而我们可以进行如下MLA操作。

\Needspace{5\baselineskip}
\hypertarget{ux4e8cux7b2cux4e00ux6b65ux751fux6210-c_kv}{%
\subsection{\texorpdfstring{二、第一步：生成 \(C_{KV}\)}{二、第一步：生成 C\_\{KV\}}}\label{ux4e8cux7b2cux4e00ux6b65ux751fux6210-c_kv}}

\Needspace{5\baselineskip}
在运用 KV Cache 的标准多头注意力中，我们每次读入一个 token（\(x_t\)），都会生成 \(d_{\mathrm{model}}\) 维（维度较高，如 1024 维）的向量 \(k_t\) 和 \(v_t\)（其中包含了各个头，如第 1-32 维为第 1 个头，以此类推），存入显存：

\Needspace{5\baselineskip}
在标准 MHA 中，当前 token 的 Key 和 Value 由输入向量直接投影得到：

\[
k_t = x_t W_K,\qquad v_t = x_t W_V
\]

\Needspace{5\baselineskip}
MLA 认为，由于不同头的向量相关性较高，这里很多维事实上是``浪费''的，我们只需要用一个共用的降维矩阵 \(W_{DKV}\) 将真正有用的信息存入一个极低维度（如 64 维）、所有头共用的向量 \(c_{KV,t}\)，后续每个头再乘以升维矩阵的不同列（保留不同头的多样性）即可。得到 \(c_{KV,t}\) 的表达式：

\[
c_{KV,t} = x_t W_{DKV}
\]

\Needspace{5\baselineskip}
位置信息单独储存：

\[
k_t^R = \mathrm{RoPE}(x_t W_{KR})
\]

\Needspace{5\baselineskip}
\hypertarget{ux4e09ux7b2cux4e8cux6b65ux8fd0ux7528ux7ed3ux5408ux5f8bux5bf9ux6ce8ux610fux529bux8ba1ux7b97ux8fdbux884cux4f18ux5316}{%
\subsection{三、第二步：运用结合律对注意力计算进行优化}\label{ux4e09ux7b2cux4e8cux6b65ux8fd0ux7528ux7ed3ux5408ux5f8bux5bf9ux6ce8ux610fux529bux8ba1ux7b97ux8fdbux884cux4f18ux5316}}

\Needspace{5\baselineskip}
原本的注意力分数可以看成 Query 与升维解压后的 Key 做点积：

\[
\mathrm{Score} \propto q_t \cdot (c_{KV} \cdot W_{UK})^T
\]

\Needspace{5\baselineskip}
利用矩阵乘法结合律 \((AB)C=A(BC)\)，可以把 Key 侧的升维矩阵 \(W_{UK}\) 移动到 Query 侧：

\[
\mathrm{Score} \propto (q_t \cdot W_{UK}^T) \cdot c_{KV}^T
\]

\Needspace{5\baselineskip}
我们先让 Query 去乘固定的参数矩阵 \(W_{UK}^T\)，得到一个新的、被变换过的 Query，记为 \(Q_{\mathrm{absorbed}}\)，则对于第 \(i\) 个注意力头，当前（第 \(t\) 个）Token 和第 \(j\) 个 Token 间的注意力权重为：

\[
\mathrm{Score}_{t,j}^{(i)}
=
\frac{Q_{\mathrm{absorbed}}^{(i)} \cdot (c_{KV,j})^T}{\sqrt{d}}
\]

由于 \(Q_{\mathrm{absorbed}}\) 是对于当前 Token 而言的，因此这一步只涉及 1 个 Token 的计算，开销极小。接下来乘以 \(C_{KV}^T\)，我们只需要存储 \(t*64\) 维的压缩向量 \(C_{KV}\)，而永远不需要在显存里复原巨大的 \(t*1024\) 维的 K 矩阵。在完成``矩阵吸收''后，这步就变成了 \(Q_{\mathrm{absorbed}}^{(i)}\) 对 \(C_{KV}^T\) 进行类似 MQA 的操作，但保留了 MHA 中不同头的表达能力。

\Needspace{5\baselineskip}
我们还需要加上解耦位置编码，单独计算 RoPE 部分的分数。从而对于第 \(i\) 个注意力头，当前（第 \(t\) 个）Token 和第 \(j\) 个 Token 间的注意力权重为：

\[
\mathrm{Score}_{t,j}^{(i)}
=
\frac{
\left(q_t^{C,(i)} \cdot (W_{UK}^{(i)})^T\right) \cdot c_{KV,j}
+
q_t^{R,(i)} \cdot (k_j^R)^T
}{\sqrt{d}}
\]

其中，第一项 \(\left(q_t^{C,(i)} \cdot (W_{UK}^{(i)})^T\right) \cdot c_{KV,j}\) 是内容分数，用于比较当前 token 与历史 token 的语义内容；第二项 \(q_t^{R,(i)} \cdot (k_j^R)^T\) 是 RoPE 位置分数，用于比较二者的相对位置信息。

\(R_{lt}^{(i)}\) 是第 \(i\) 个头独有的、携带位置信息的 Query 向量，可理解为上式中的 \(q_t^{R,(i)}\)；\(k_j^R\) 通常是各头共享的 RoPE Key 向量。它们做点积后，相当于在标准注意力中保留位置匹配能力，使模型仍然能判断 token 之间的相对距离。

\Needspace{5\baselineskip}
\hypertarget{ux56dbux8fd0ux7528ux5206ux914dux5f8bux5bf9ux4fe1ux606fux805aux5408ux6ce8ux610fux529bux6743ux91cdvux8fdbux884cux4f18ux5316}{%
\subsection{四、运用分配律对信息聚合（注意力权重*V）进行优化}\label{ux56dbux8fd0ux7528ux5206ux914dux5f8bux5bf9ux4fe1ux606fux805aux5408ux6ce8ux610fux529bux6743ux91cdvux8fdbux884cux4f18ux5316}}

\Needspace{5\baselineskip}
设 \(c_{KV,j}\) 向量的维度为 \(d_c\)，升维后每个头的特征维度为 \(d_h\)。在第 \(i\) 个头中，注意力分布、压缩潜变量和 Value 升维矩阵分别为：

\[
P_{t,:}^{(i)} = \mathrm{Softmax}(\mathrm{Score}_{t,:}^{(i)}),\qquad
P_{t,j}^{(i)} =
\frac{\exp(\mathrm{Score}_{t,j}^{(i)})}{\sum_{r=1}^{t}\exp(\mathrm{Score}_{t,r}^{(i)})},\qquad
c_{KV,j}\in \mathbb{R}^{d_c},\qquad
W_{UV}^{(i)}\in \mathbb{R}^{d_c \times d_h}
\]

\Needspace{5\baselineskip}
如果按标准路径先把每个历史 token 的压缩潜变量升维成 Value，再做加权求和，则第 \(i\) 个头的输出是：

\[
y_t^{(i)}
=
\sum_{j=1}^{t}
\left(P_{t,j}^{(i)} \cdot (c_{KV,j} \cdot W_{UV}^{(i)})\right)
\]

这种写法对每个历史位置 \(j\) 都要执行一次 \(c_{KV,j} \cdot W_{UV}^{(i)}\) 的升维计算，然后再乘以标量权重 \(P_{t,j}^{(i)}\) 并求和。单个头的主要复杂度约为 \(O(t \cdot d_c \cdot d_h)\)。

\Needspace{5\baselineskip}
DeepSeek 利用矩阵乘法的分配律，把固定的升维矩阵 \(W_{UV}^{(i)}\) 移到求和符号之外：

\[
\begin{aligned}
y_t^{(i)}
&= \sum_{j=1}^{t}\left(P_{t,j}^{(i)} \cdot c_{KV,j} \cdot W_{UV}^{(i)}\right) \\
&= \left(\sum_{j=1}^{t}P_{t,j}^{(i)} \cdot c_{KV,j}\right) \cdot W_{UV}^{(i)}
\end{aligned}
\]

\Needspace{5\baselineskip}
于是计算可以拆成两个阶段。第一阶段先在压缩空间内聚合：

\[
\tilde c_{sum}^{(i)}
=
\sum_{j=1}^{t}\left(P_{t,j}^{(i)} \cdot c_{KV,j}\right)
\]

\(\tilde c_{sum}^{(i)}\) 是一个 \(d_c\) 维低维向量。虽然 \(c_{KV,j}\) 来自共享 KV Cache，但每个头的注意力权重 \(P_{t,j}^{(i)}\) 不同，所以每个头聚合出的 \(\tilde c_{sum}^{(i)}\) 也不同。这个阶段只做标量乘低维向量再求和，复杂度约为 \(O(t \cdot d_c)\)。

\Needspace{5\baselineskip}
第二阶段再做延迟升维：

\[
y_t^{(i)} = \tilde c_{sum}^{(i)} \cdot W_{UV}^{(i)}
\]

这一步把聚合后的低维结果映射回第 \(i\) 个头的 \(d_h\) 维特征空间。由于升维只对聚合后的结果做一次，而不是对 \(t\) 个历史 token 分别做 \(t\) 次，单个头这部分复杂度约为 \(O(d_c \cdot d_h)\)。

\Needspace{5\baselineskip}
计算出所有 \(H\) 个头的 \(y_t^{(i)}\) 后，将它们拼接起来，并通过输出投影矩阵 \(W_O\) 得到最终的隐状态输出：

\[
y_t
=
\mathrm{Concat}\left(y_t^{(1)}, y_t^{(2)}, \ldots, y_t^{(H)}\right) \cdot W_O
\]

总体而言，MLA 把``先逐 token 升维再聚合''改成``先在压缩空间聚合再延迟升维''，主要复杂度从 \(O(t \cdot d_c \cdot d_h)\) 变为 \(O(t \cdot d_c + d_c \cdot d_h)\)，当历史长度 \(t\) 较大时可近似理解为降低约 \(d_h\) 倍。

\FloatBarrier
\Needspace{5\baselineskip}
\hypertarget{ux53c2ux8003ux6587ux732e-95}{%
\subsection{参考文献}\label{ux53c2ux8003ux6587ux732e-95}}

\vspace{0.25em}
\begin{itemize}
\tightlist
\item
  DeepSeek-AI. (2024). \href{https://arxiv.org/abs/2405.04434}{DeepSeek-V2: A Strong, Economical, and Efficient Mixture-of-Experts Language Model}. arXiv:2405.04434.
\end{itemize}

\clearpage
